%======================================================================
\chapter{Moore's axioms and the bridge to topology}\label{ch:axioms}
%======================================================================

To prove that the space of elements of a Moore decomposition is a plane, Moore uses his own
axiomatic characterization of the plane (1916, 1919): a set of ``points'' with a class of subsets
called ``regions'' satisfying eight axioms is a plane. Chapters~\ref{ch:axioms}--\ref{ch:plane}
prove this characterization; Chapter~\ref{ch:quotient} verifies the axioms for the space of elements.

This chapter states the axioms (Moore 1916, \S2) and shows that Moore's notions agree with those of
the topology generated by the regions (the ``bridge''). Then the topology does the work: the space is
metrizable, closures of regions are compact, two points of a connected open set are joined by an arc
(Theorem~15, by Moore's nested chains), and Moore's arcs are the topological ones. We cite Moore's
paper \emph{On the foundations of plane analysis situs} (1916) as F.A., as Moore does.

\section{Points, regions and the axioms}\label{sec:axioms}

A \emph{Moore space} is a set $S$ together with a class of subsets of $S$, the \emph{regions}.

\begin{definition}[Moore's derived notions]\label{def:moore-notions}
Let $M\subseteq S$.
\begin{itemize}
\item $P$ is a \emph{limit point} of $M$ if every region containing $P$ contains a point of $M$
other than $P$.
\item The \emph{boundary} $\bd M$ is the set of points $X$ such that every region containing $X$
contains a point of $M$ and a point not in $M$. $M'=M\cup\bd M$, and the \emph{exterior} of $M$ is
$S\setminus M'$.
\item $M$ is \emph{connected} if, however $M$ is divided into two disjoint nonempty sets, one of them
contains a limit point of the other. $M$ is \emph{closed} if it contains all its limit points.
\item $M$ is \emph{bounded} if $M\subseteq(R_1\cup\dots\cup R_n)'$ for finitely many regions; $M$ has
the \emph{Heine--Borel property} if every cover of $M$ by regions has a finite subcover.
\item A \emph{domain} is a connected set each point of which lies in a region contained in the set.
\item An \emph{arc from $A$ to $B$} ($A\ne B$) is a bounded, closed, connected set containing $A$ and
$B$ with no connected proper subset containing both. A \emph{simple closed curve} is the union of two
arcs from $A$ to $B$ with only $A$ and $B$ in common.
\item A \emph{simple chain from $A$ to $B$} is a finite sequence of regions $R_0,\dots,R_{n-1}$ such
that $A$ lies in $R_0$ only, $B$ in $R_{n-1}$ only, and for $i<j$, $R_i\cap R_j\ne\emptyset$ if and only
if $j=i+1$.
\end{itemize}
\end{definition}

\formal{\lean{HPC.Moore.Space}, \lean{HPC.Moore.Space.IsLimitPt}, \lean{HPC.Moore.Space.bdry},
\lean{HPC.Moore.Space.cl}, \lean{HPC.Moore.Space.ext}, \lean{HPC.Moore.Space.IsConn},
\lean{HPC.Moore.Space.IsBoundedM}, \lean{HPC.Moore.Space.IsDomain}, \lean{HPC.Moore.Space.IsArc},
\lean{HPC.Moore.Space.IsSCC}, \lean{HPC.Moore.Space.IsSimpleChain}.}

In Moore's definition of ``connected'' the two parts must be nonempty (otherwise no set would be
connected); so read, the empty set is connected, i.e.\ Moore's ``connected'' is ``preconnected''.

\begin{definition}[Moore's axioms, F.A.\ pp.~132--133]\label{def:moore-axioms}
\leavevmode
\begin{description}
\item[Axiom 1.] There is a sequence of regions $K_1,K_2,\dots$ such that (1) every point lies in $K_n$
for arbitrarily large $n$; (2) if $P\ne\bar P$ are points of a region $R$, then for all large $n$,
$P\in K_n$ implies $K_n'\subseteq R\setminus\set{\bar P}$.
\item[Axiom 2.] Every region is connected.
\item[Axiom 3.] The exterior of every region is connected.
\item[Axiom 4.] For every region $R$, $R'$ has the Heine--Borel property.
\item[Axiom 5.] There is an infinite set of points with no limit point.
\item[Axiom 6.] If $R,\bar R$ are regions and $P\in\bar R$ is a boundary point of $R$, there are regions
$K,\bar K\subseteq\bar R$ with $P\in\bar K$, $K\subseteq R$, and every boundary point of $R$ in $\bar K$ a
boundary point of $K$.
\item[Axiom 7.] The same as Axiom~6 with ``$L$ in the exterior of $R$'' in place of ``$K\subseteq R$''.
\item[Axiom 8.] Every simple closed curve is the boundary of some region.
\end{description}
Axiom~$5_1$ (a footnote of F.A.) says that there are at least three points; it follows from Axiom~5.
\end{definition}

\formal{\lean{HPC.Moore.Space.Axiom1} \dots\ \lean{HPC.Moore.Space.Axiom8},
\leansub{HPC.Moore.Space.Axiom5}{1}, \leansub{HPC.Moore.axiom5}{1}\lean{_of_axiom5}.}

The model to keep in mind: $S=\R^2$, the regions the interiors of simple closed curves (or the open
squares). The \emph{fundamental regions} $K_n$ of Axiom~1 play the role of a countable basis of
shrinking sets.

\section{The bridge}

Let $\tau$ be the topology on $S$ generated by the regions (the smallest topology in which every region
is open).

\begin{proposition}[The bridge]\label{prop:bridge}
Assume Axiom~1.
\begin{enumerate}[label=(\alph*)]
\item The regions form a basis of $\tau$: if $P$ lies in two regions, some region about $P$ lies in
both (F.A.\ Theorem~6).
\item $P$ is a limit point of $M$ in Moore's sense if and only if it is an accumulation point of $M$ in
$\tau$. Consequently $M'$ is the closure of $M$, Moore's closed sets are the closed sets, $\bd M$ is the
frontier of $M$ minus the interior points, the exterior of a region is the complement of its closure,
Moore's connected sets are the preconnected sets, and the domains are the open connected sets.
\item The Heine--Borel property is compactness.
\item $\tau$ is $T_1$ and regular.
\item With Axioms~3 and $5_1$: every nonempty region is infinite (F.A.\ Theorem~4), the fundamental regions
form a countable basis, and $\tau$ is metrizable.
\item With Axiom~4: closures of regions are compact; a set is bounded if and only if it lies in a
compact set.
\end{enumerate}
\end{proposition}

\begin{proof}
(a) Let $P\in K\cap L$. If $K$ has a point $\bar P\ne P$, Axiom~1 gives $n$ with $P\in K_n$ and
$K_n'\subseteq K\setminus\set{\bar P}$; applying Axiom~1 again inside $L$ (if $L\ne\set P$) and taking $n$
large enough for both, $K_n\subseteq K\cap L$. If $K=\set P$ or $L=\set P$, that region itself will do.

(b) With (a), ``every region containing $P$'' means ``every basic neighbourhood of $P$''. The other
statements are rewritings; for example, $M$ is connected in Moore's sense iff in every partition into
two nonempty pieces one meets the closure of the other, i.e.\ iff $M$ is preconnected.

(c) Covers by regions refine arbitrary open covers, since regions form a basis.

(d) Points are closed: if $\bar P\ne P$, Axiom~1 (2) applied to a region containing both (if there is
none, any region about $P$ misses $\bar P$) gives a region about $P$ missing $\bar P$. Regularity:
given $P$ in a region $R$, if $R$ has a second point, Axiom~1 (2) gives $K_n\ni P$ with $K_n'\subseteq R$.

(e) Let $O\in R$. By Axiom~$5_1$ there are points $P_1\ne P_2$ different from $O$. By (d) there is a
region $R_1\ni P_1$ missing $O$ and $P_2$, and a region $K\ni P_1$ with $K'\subseteq R_1$. Then $O$ and $P_2$
lie in the exterior of $K$, which is connected (Axiom~3). Divide it into $\set O$ and the rest (which
contains $P_2$): one part contains a limit point of the other. A point is never a limit point of the
single point $O$ (by (d)), so $O$ is a limit point of the exterior of $K$. Hence every region about $O$, in
particular $R$, contains infinitely many points (a region about a limit point of $M$ contains infinitely
many points of $M$, by (d) and Axiom~1). So there are no isolated points; every region about $P$ has a
second point, Axiom~1 (2) applies, and the $K_n$ containing $P$ with large $n$ lie in any given region
about $P$: the $K_n$ form a countable basis. A regular second countable $T_1$ space is metrizable.

(f) Axiom~4 and (c). A bounded set lies in the compact set $(R_1\cup\dots\cup R_n)'=\cl{R_1}\cup\dots\cup
\cl{R_n}$; conversely a compact set is covered by finitely many regions, so it is bounded.
\end{proof}

\formal{\lean{HPC.Moore.Space.topology}, \lean{HPC.Moore.fa_theorem6},
\lean{HPC.Moore.isTopologicalBasis_region}, \lean{HPC.Moore.isLimitPt_iff},
\lean{HPC.Moore.cl_eq_closure}, \lean{HPC.Moore.isClosedM_iff}, \lean{HPC.Moore.isConn_iff},
\lean{HPC.Moore.isDomain_iff}, \lean{HPC.Moore.heineBorel_iff}, \lean{HPC.Moore.t1Space},
\lean{HPC.Moore.regularSpace}, \lean{HPC.Moore.fa_theorem4}, \lean{HPC.Moore.secondCountable},
\lean{HPC.Moore.metrizable}, \lean{HPC.Moore.isCompact_closure_region},
\lean{HPC.Moore.isBoundedM_of_isCompact}.}

What matters for the rest of the volume is the conclusion: \emph{under Axioms~1, 3, 4
and $5_1$, $S$ is a metrizable space in which the regions form a basis of open sets with compact
closures.}

\begin{remark}\label{rem:empty-region}
Moore's Theorem~4 (``every region contains infinitely many points'') and Theorem~19 (``every region has
a boundary point'') are printed for all regions, and his proofs begin ``there exists in $R$ at least one
point''. The axioms do not exclude $\emptyset$ as a region: adding $\emptyset$ to the regions changes none of
the derived notions and preserves Axioms~1--8. So the printed statements fail in any model with
$\emptyset$ added, and ``nonempty'' must be added (\lean{HPC.Moore.not_fa_theorem4_printed},
\lean{HPC.Moore.not_fa_theorem19_printed}). Similarly Theorem~14 needs every $M_n$ nonempty
(\lean{HPC.Moore.not_fa_theorem14_printed}).
\end{remark}

\section{Elementary consequences}

The following theorems of F.A.\ \S\S3--5 are, after the bridge, standard facts about a metrizable space
with a basis of connected open sets with compact closures. We list them with the topological reason.

\begin{theorem}[F.A.\ Theorems~1--3, 5, 7--9, 11--14, 17--19, 22]\label{thm:fa-elementary}
Assume Axioms~1--5.
\begin{enumerate}[label=(\alph*)]
\item (1) No point of a region is a boundary point of it. (2) Every region about a limit point of $M$
contains infinitely many points of $M$. (3) A point not in a finite set $M$ lies in a region missing
$M$. (7) A limit point of $M\cup N$ is a limit point of $M$ or of $N$. (8), (9) Limit points are limits
of sequences, and a convergent sequence has no other limit point.
\item (11), (12) Closed bounded sets are compact. (13) Every infinite bounded set has a limit point.
(14) A decreasing sequence of nonempty bounded sets $M_n\supseteq M_{n+1}'$ has nonempty closed
intersection.
\item (17), (18) The exterior of a region contains an infinite set without limit point and is unbounded.
(19) Every nonempty region has a boundary point. (22) The whole space is connected.
\end{enumerate}
\end{theorem}

\begin{proof}
(a) and (b) follow from the bridge (Proposition~\ref{prop:bridge}) and metrizability, and (14) is the
nested intersection property of compact sets.
(17) The infinite set $M$ of Axiom~5 meets the bounded set $R'$ in a finite set (by (13)), so $M\setminus
R'$ will do; (18) follows by (13).
(22) If $S=U\cup V$ with $U,V$ disjoint, closed and nonempty, take a region $R\subseteq U$ (a region about a
point of $U$, inside the open set $U$). Its exterior is connected (Axiom~3) and contains $V$, and so lies in
$V$; hence $U\subseteq R'$, which is compact (Axiom~4). Likewise $V$ is compact, so $S$ is compact, and
then every infinite set has a limit point, against Axiom~5.
(19) Let $P\in R$; by (22) $R$ is not closed and open unless $R=S$, which is excluded by (18); so $\fr R$ is
nonempty, and as $R$ is open, $\fr R=\bd R$.
\end{proof}

\formal{\lean{HPC.Moore.fa_theorem1}, \lean{HPC.Moore.fa_theorem2}, \lean{HPC.Moore.fa_theorem3},
\lean{HPC.Moore.fa_theorem5}, \lean{HPC.Moore.fa_theorem7}, \lean{HPC.Moore.fa_theorem8},
\lean{HPC.Moore.fa_theorem9}, \lean{HPC.Moore.fa_theorem11}, \lean{HPC.Moore.fa_theorem12},
\lean{HPC.Moore.fa_theorem13}, \lean{HPC.Moore.fa_theorem14}, \lean{HPC.Moore.fa_theorem17},
\lean{HPC.Moore.fa_theorem18}, \lean{HPC.Moore.fa_theorem19}, \lean{HPC.Moore.isConn_univ},
\lean{HPC.Moore.fa_theorem22}.}

Each formal statement carries only the axioms its proof uses; for example Theorems~11--14 use only
Axioms~1 and~4.

\section{Simple chains and Theorem 15}

\begin{theorem}[F.A.\ Theorem~10]\label{thm:fa10}
Let $M$ be connected, $A\ne B$ points of $M$, and $\mathcal G$ a set of regions covering $M$. There is a
simple chain from $A$ to $B$ all of whose links belong to $\mathcal G$.
\end{theorem}

\begin{proof}
Call a \emph{weak chain} from $A$ to $B$ a finite sequence of sets with $A$ in the first, $B$ in the last,
and consecutive members meeting. Let $S_A$ be the set of points $X$ of $M$ such that there is a weak chain
of members of $\mathcal G$ from $A$ to $X$. Then $S_A$ and $M\setminus S_A$ are both relatively open in $M$
(if $X\in M$ lies in $R\in\mathcal G$, every point of $M\cap R$ is in $S_A$ iff $X$ is: append $R$). So
neither contains a limit point of the other, and since $M$ is connected and $A\in S_A$, $S_A=M\ni B$.

Now take a weak chain from $A$ to $B$ of members of $\mathcal G$ of minimal length; it is a simple chain.
Indeed, if $A$ lay in a link $R_i$ with $i>0$, the chain $R_i,\dots,R_{n-1}$ would be shorter; similarly
for $B$; and if $R_i\cap R_j\ne\emptyset$ with $j>i+1$, then $R_0,\dots,R_i,R_j,\dots,R_{n-1}$ would be a
shorter weak chain. Note that the minimal chain is a \emph{sublist} of the original, with the order of the
links preserved.
\end{proof}

\formal{\lean{HPC.Moore.fa_theorem10}, \lean{HPC.Moore.exists_simpleChain_sublist}.}

\begin{theorem}[F.A.\ Theorem~15]\label{thm:fa15}
Assume Axioms~1--4 and~$5_1$. Two distinct points of a domain $M$ are joined by an arc lying in $M$.
\end{theorem}

\begin{proof}
By the bridge, $S$ is metrizable (fix a metric $d$), $M$ is open and connected, and the regions are
connected open sets with compact closures, with small regions about every point inside every open set.

\emph{The chains.} We construct simple chains $C_1,C_2,\dots$ of regions from $A$ to $B$ such that
\begin{enumerate}[label=(\arabic*)]
\item each link of $C_1$ has its closure in $M$; each link of $C_{n+1}$ has its closure in a single link
of $C_n$, its \emph{owner};
\item the owner map is monotone: if a link precedes another in $C_{n+1}$, the owner of the first is not
after the owner of the second;
\item every link of $C_n$ lies in a ball of radius $1/n$.
\end{enumerate}
$C_1$: apply Theorem~\ref{thm:fa10} to $M$ and the regions of diameter $<1$ with closure in $M$. Given
$C_n=(R_0,\dots,R_{k-1})$, choose points $p_i\in R_{i-1}\cap R_i$ ($0<i<k$), $p_0=A$, $p_k=B$. Inside the
connected open set $R_i$, Theorem~\ref{thm:fa10} gives a simple chain from $p_i$ to $p_{i+1}$ of regions of
radius $<1/(n+1)$ with closures in $R_i$ (if $p_i=p_{i+1}$, a single such region). Laying these chains end
to end gives a weak chain from $A$ to $B$, in which the links coming from $R_i$ come before those coming
from $R_{i+1}$; its shortest sub-chain (proof of Theorem~\ref{thm:fa10}) is a sublist, so (1)--(3) hold
with the owner of a link the region $R_i$ it came from.

\emph{The continuum.} Let $F_n$ be the union of the closures of the links of $C_n$. Each $F_n$ is compact
and connected (closures of links are connected and consecutive ones meet), $F_{n+1}\subseteq F_n$ by (1),
and $A,B\in F_n$. So $C=\bigcap_nF_n$ is a continuum containing $A$ and $B$, and $C\subseteq F_1\subseteq M$.

\emph{The order.} A point lies in at most two links of a simple chain, and if in two, they are
consecutive. For $X,Y\in C$ write $X\prec Y$ if at some stage $n$ every link of $C_n$ containing $X$
precedes every link containing $Y$. By (2) this persists to all later stages. If $X\ne Y$, then at a stage
with $2/n<d(X,Y)/2$ no link containing $X$ meets a link containing $Y$, so the (one or two, consecutive)
links containing $X$ all precede, or all follow, those containing $Y$: $X\prec Y$ or $Y\prec X$. Also $A\prec
P$ and $P\prec B$ for $P\ne A,B$ ($A$ lies in the first link only, $B$ in the last only).

\emph{Irreducibility.} Fix $P\in C\setminus\set{A,B}$ and let $S_A=\set{X\in C:X\prec P}$,
$S_B=\set{X\in C:P\prec X}$. These are disjoint and cover $C\setminus\set P$, and $A\in S_A$, $B\in S_B$. We
show $S_A$ is open in $C$ (likewise $S_B$). Let $X\in S_A$; choose $n$ so large that $X\prec P$ at stage $n$
and no two meeting links of $C_n$ contain $X$ and $P$ respectively. Let $L$ be a link of $C_n$ containing $X$
and $Y\in C\cap L$. The links containing $Y$ are $L$ or adjacent to $L$; the links containing $P$ follow $L$
and are not adjacent to it. So every link containing $Y$ precedes every link containing $P$, $Y\prec P$, and
$C\cap L\subseteq S_A$ is a neighbourhood of $X$ in $C$. Hence $C\setminus\set P$ is the union of two disjoint
relatively open sets, one containing $A$ and the other $B$, and every connected subset of $C$ containing $A$
and $B$ contains $P$. So $C$ is irreducible between $A$ and $B$, and it is an arc by
Corollary~\ref{cor:arc-char}.
\end{proof}

\formal{\lean{HPC.Moore.exists_chains}, \lean{HPC.Moore.exists_irreducible},
\lean{HPC.Moore.fa_theorem15}.}

\begin{corollary}[F.A.\ Theorem~16]\label{cor:fa16}
Two points of a region are joined by an arc in the region; two points of the exterior of a region are
joined by an arc in the exterior.
\end{corollary}

\begin{proof}
A region is a domain (Axiom~2); its exterior is open (the closure is closed) and connected (Axiom~3).
\end{proof}

\formal{\lean{HPC.Moore.fa_theorem16}.}

\section{Moore's arcs and simple closed curves}

\begin{proposition}\label{prop:moore-arcs}
Assume Axioms~1, 3, 4 and $5_1$. A set is an arc from $A$ to $B$ in Moore's sense if and only if it is an
arc from $A$ to $B$ in the topology $\tau$ (Definition~\ref{def:arc}); and it is a simple closed curve in
Moore's sense if and only if it is one in $\tau$.
\end{proposition}

\begin{proof}
Bounded and closed means compact (Proposition~\ref{prop:bridge}), connected means connected, and
Corollary~\ref{cor:arc-char} applies in the metrizable space $(S,\tau)$.
\end{proof}

\formal{\lean{HPC.Moore.isArc_iff}, \lean{HPC.Moore.isSCC_iff_tscc}.}

From now on all notions are topological: ``region'' means a member of the given basis, ``arc'' and
``simple closed curve'' have their topological meaning, and $\cl M$, $\fr M$ are closure and frontier.
Moore's unnumbered facts about arcs (F.A.\ pp.~139--140: splitting an arc at a point, first and last points
on a closed set, the two arcs of a simple closed curve) are the results of Chapter~\ref{ch:arcs}
(\lean{HPC.Moore.fa_sec5_arc_split}, \lean{HPC.Moore.fa_sec5_first_point}, \lean{HPC.Moore.fa_scc_two_arcs}).

\section*{Exercises}
\addcontentsline{toc}{section}{Exercises}

\begin{exercise}
Check that $\R^2$ with the open squares as regions satisfies Axioms~1--5. (Axioms~6--8 for this choice of
regions are not obvious: Axiom~8 says that every simple closed curve bounds a square!) Show that with the
interiors of simple closed curves as regions all eight axioms hold, using the Jordan curve theorem and the
Schoenflies theorem.
\end{exercise}

\begin{exercise}
Show that the sphere $S^2$ with the interiors of simple closed curves (either complementary domain) as
regions satisfies Axioms~1--4 and 6--8 but not Axiom~5. Which step of the proof of Theorem~\ref{thm:fa-elementary}(22)
uses Axiom~5?
\end{exercise}

\begin{exercise}
Show that the construction in the proof of Theorem~\ref{thm:fa15} works in any connected, locally connected,
locally compact metric space, and deduce that such a space is arcwise connected.
\end{exercise}
